
%%%%%%%%%%%%%%%%%%%%%%%%%%%%%%%%%%%%%%%%%%%%%%%%%%%%%%%%%%%%%%%%%%%%%%
% How to use writeLaTeX: 
%
% You edit the source code here on the left, and the preview on the
% right shows you the result within a few seconds.
%
% Bookmark this page and share the URL with your co-authors. They can
% edit at the same time!
%
% You can upload figures, bibliographies, custom classes and
% styles using the files menu.
%
%%%%%%%%%%%%%%%%%%%%%%%%%%%%%%%%%%%%%%%%%%%%%%%%%%%%%%%%%%%%%%%%%%%%%%

\documentclass[12pt]{article}

\usepackage{sbc-template}

\usepackage{graphicx,url}

\usepackage[brazilian]{babel}   
\usepackage[T1]{fontenc} % Codificação para fontes com suporte a caracteres acentuados
\usepackage[utf8]{inputenc} % Codificação UTF-8 (necessário apenas para pdflatex)
\usepackage{multicol}
\usepackage{multirow}
\usepackage{graphicx,url}
\usepackage{amsmath,amssymb,amsfonts}
\usepackage{algorithmic}
\usepackage{textcomp}
\usepackage{url}
\usepackage{enumerate}
\usepackage{multirow}
\usepackage{soul}
\usepackage{verbatim}
\usepackage[table,xcdraw]{xcolor}
\usepackage[shortlabels]{enumitem}
\usepackage{scalefnt} % pacote para diminuir fonte das tabelas
\usepackage{comment} % pacote de comentarios
\usepackage{hanging} % recuo invertido nas referências
\usepackage{lscape}
\usepackage{fancyhdr}
\usepackage{sbgames}
\usepackage[super]{nth}

%previne hifenização
\tolerance=1
\emergencystretch=\maxdimen
\hyphenpenalty=10000
\hbadness=10000

\sloppy

\newcommand{\yearconference}{2025}
\newcommand{\edition}{XXIV}
\newcommand{\city}{Recife -- PE}
\newcommand{\period}{October 16th - 29th}



\title{Mitigação de Viés Algorítmico em Cenários de Desbalanceamento Sociodemográfico: Uma Análise Comparativa entre Foundation Models (TabPFN) e Gradient Boosting (XGBoost) nas Eleições Municipais do Ceará \\
\medskip
\small{
    \textit{Algorithmic Bias Mitigation in Sociodemographic Imbalance Scenarios: A Comparative Analysis of Foundation Models (TabPFN) and Gradient Boosting (XGBoost) in Ceará Municipal Elections}
}
}

\author{Luana Teles Alves\inst{1}, Luann Alves Pereira de Lima\inst{1}, Thiago Nerton Macedo Alves\inst{1},\\ Vitória do Nascimento Pontes\inst{1}}

\address{Centro de Ciências e Tecnologia (CCT) -- Universidade Federal do Cariri (UFCA)\\
  Juazeiro do Norte -- CE -- Brasil
  \email{luana.teles@aluno.ufca.edu.br, luann.alves@aluno.ufca.edu.br}
  \email{thiago.nerton@aluno.ufca.edu.br, vitoria.pontes@aluno.ufca.edu.br}
}



\begin{document} 

\maketitle

\thispagestyle{plain}


\begin{abstract}
  Under cross-validation ($K=10$, $N=2,178$) with demographic attributes, both
  achieved equivalent predictive performance (MCC $\approx$ 0.28). However, XGBoost
  exhibited significant gender and racial bias (DIR = 1.17 and 0.88, respectively),
  amplifying historical patterns of self-selection. In contrast, TabPFN maintained
  relative parity in both dimensions (DIR $\approx$ 1.01 and 0.91), evidencing that
  architectural differences between foundation models and gradient boosting can produce
  relevant fairness disparities even under equivalent predictive accuracy.
\end{abstract}
     
\begin{resumo} 
  Sob validação cruzada ($K=10$, $N=2.178$) com atributos demográficos, ambos
  atingiram desempenho preditivo equivalente (MCC $\approx$ 0,28). Contudo, o
  XGBoost apresentou viés de gênero e raça significativo (DIR = 1,17 e 0,88,
  respectivamente), amplificando padrões históricos de auto-seleção. Em
  contrapartida, o TabPFN manteve paridade relativa em ambas as dimensões
  (DIR $\approx$ 1,01 e 0,91), evidenciando que diferenças arquiteturais entre
  foundation models e gradient boosting podem gerar disparidades de equidade
  relevantes mesmo sob acurácia preditiva equivalente.
\end{resumo}

\section{Introdução}

A predição de resultados eleitorais em eleições majoritárias no Ceará enfrenta desafios
metodológicos complexos dada a natureza desequilibrada dos dados – muitos candidatos para
poucas vagas – e a influência de barreiras históricas sociodemográficas. O cargo de
prefeito, em particular, exige uma análise que vá além da simples performance estatística,
visando na equidade e integridade dos modelos de decisão.

Com o  uso dos algoritmos de aprendizado de máquina cada vez mais presente nos processos de 
tomadas de decisões em diferentes domínios sociais, surge a preocupação urgente em relação a
amplificação de preconceitos contra um grupo de pessoas, baseados em características sensíveis
como raça e gênero. Modelos baseados em árvores de decisão, frequentemente utilizados em
dados tabulares, operam sob a lógica de otimização da acurácia global, o que pode resultar
em padrões de classificação distintos para diferentes grupos populacionais [1]. 

Historicamente algoritmos que se baseiam no Gradient Boosting, como o XGBoost, viraram padrão
na indústria quando se trata de dados tabulares devido a suas taxas de acurácia geralmente mais
altas. Sua eficácia reside na construção iterativa de árvores de decisão, que, através da
otimização de uma função de perda global, conseguem capturar relações não-lineares complexas
com alta eficiência computacional [3]. No entanto, esse paradigma de "minimização de perda empírica"
impõe limitações estruturais: modelos dependentes exclusivamente de conjuntos de dados locais são
suscetíveis à cristalização de vieses presentes no histórico de treinamento, além de exigirem uma
otimização exaustiva de hiperparâmetros para alcançar estabilidade preditiva em “small data”.

Nesse cenário, emerge o paradigma dos Foundation Models para dados tabulares, representados de forma
proeminente pelo TabPFN (Prior-Data Fitted Network). Diferente dos métodos de boosting, o TabPFN é
uma rede neural baseada na arquitetura Transformer que realiza inferência através de Aprendizado em
Contexto (In-Context Learning) [4, 8]. Em vez de aprender os pesos durante um processo de treinamento
interativo, o modelo utiliza um conhecimento prévio consolidado durante um vasto treinamento em datasets
sintéticos e reais, permitindo que a predição em novos dados seja executada como uma consulta bayesiana
sobre as distribuições previamente aprendidas [4, 9]. Embora a literatura internacional apresenta avanços
significativos na mitigação desses vieses em aprendizado de máquina, sua aplicação prática na área eleitoral
em específico é baixa ou quase nula. Estudos recentes têm comparado o desempenho de arquiteturas de fundação
e métodos tradicionais em benchmarks sintéticos [2, 4], porém, carecem de validação em contextos reais de
eleições municipais brasileiras. 

Neste artigo, o objetivo central é realizar uma análise comparativa do desempenho preditivo e das métricas
de equidade entre o TabPFN e o XGBoost, utilizando dados do Tribunal Superior
Eleitoral [5]. A investigação baseia-se em uma abordagem de validação cruzada estratificada
(StratifiedKFold, $K=10$), buscando observar como essas arquiteturas se comportam frente ao desbalanceamento
das classes e à presença de atributos sensíveis.  Com base nas propriedades arquiteturais contrastantes
entre o TabPFN e o XGBoost, formulam-se as seguintes hipóteses de trabalho: H1 (Paridade de Performance):
Ambas as arquiteturas apresentam desempenho preditivo estatisticamente equivalente em cenários de small
data eleitoral, conforme mensurado pelo Coeficiente de Correlação de Matthews (MCC). H2 (Divergência de
Equidade de Gênero): O XGBoost apresenta viés de gênero mensurável ($DIR \neq 1$), enquanto o TabPFN,
devido ao seu paradigma de inferência bayesiana amortizada, não evidencia disparidade estatisticamente
significativa nesta dimensão. H3 (Divergência de Equidade Racial): O XGBoost amplifica desigualdades
raciais latentes nos dados históricos, enquanto o TabPFN atenua tais disparidades, mantendo os indicadores
de equidade próximos à neutralidade ($DIR \approx 1$). Os resultados obtidos indicam que, embora ambas as
arquiteturas alcancem desempenho preditivo estatisticamente equivalente, o XGBoost apresenta viés de
gênero e raça confirmados por indicadores de Impacto Disparado (DIR). Em contraste, o TabPFN não evidencia
disparidades significativas, sugerindo que as diferenças arquiteturais entre modelos de fundação e métodos
de Gradient Boosting podem produzir impactos distintos na equidade algorítmica, mesmo sob acurácia preditiva
equiparável O restante deste artigo está estruturado da seguinte forma: a Seção 3 detalha o protocolo
experimental, incluindo a descrição dos dados do TSE e o desenho do split temporal. A Seção 4 apresenta
os resultados quantitativos, com foco nas métricas de performance e justiça algorítmica. Por fim, a Seção
5  discute as implicações desses achados para a integridade do processo eleitoral, apresentando as
considerações finais e perspectivas para trabalhos futuros.

\section{Metodologia} \label{sec:metodologia}

\subsection{Dados}

\subsubsection{Fonte de dados}

Os dados utilizados neste estudo foram obtidos diretamente do repositório de dados
abertos do Tribunal Superior Eleitoral (TSE) [5], órgão responsável pela organização
e fiscalização do processo eleitoral brasileiro. Para fins desta pesquisa, foi feito
um recorte temporal compreende quatro ciclos eleitorais municipais, entre 2012 e 2024,
e aplicou-se um filtro geográfico restrito ao estado do Ceará. Além disso, foram mantidas
apenas as candidaturas ao cargo de prefeito municipal, o que delimita um contexto de small
data, no qual o modelo TabPFN costuma se destacar em relação ao XGBoost.

\subsubsection{Unidade de Análise}
A unidade de análise adotada é o indivíduo candidato, de modo que cada instância do conjunto
de dados corresponde a um registro único de candidatura. O conjunto de dados totaliza 2.216
instâncias, distribuídas ao longo dos quatro pleitos analisados. Para cada candidato, o
conjunto de dados contém informações referentes ao município de disputa, características
sociodemográficas (cor/raça autodeclarada, estado civil, faixa etária, gênero e grau de
instrução), dados político-partidários (partido de filiação e número de candidatura),
ocupação declarada, situação final na totalização dos votos, quantitativo de votos válidos,
votos nominais recebidos e data de carga do registro no sistema do TSE.

\subsubsection{Variáveis}

As variáveis foram organizadas em dois grupos funcionais. O primeiro reúne as características
sociodemográficas e político-institucionais do candidato que serão utilizadas como preditores:
gênero, raça/cor autodeclarada, faixa etária, estado civil, grau de instrução, ocupação declarada,
partido político de filiação, número do candidato e município. O segundo grupo diz respeito à
variável dependente, a situação da totalização indicando se aquele indivíduo foi eleito ou não.

Variáveis que introduzem data leakage foram excluídas do processo de modelagem. Em particular,
os campos de votos válidos e votos nominais foram descartados, uma vez que constituem informações
geradas pelo próprio processo eleitoral que se busca predizer. Afinal, sua inclusão como preditores
produziria modelos trivialmente precisos, porém sem qualquer capacidade de generalização para
cenários de predição pré-eleitoral (INSERIR REFERENCIA PARA DATALEAKAGE).

\subsubsection{Distribuições relevantes}

O conjunto de dados apresenta desbalanceamento de classes expressivo na variável-alvo: dos 2.216
registros, 1.411 correspondem a candidatos não eleitos (63,7%) e 767 a candidatos eleitos (34,6%)
— o restante corresponde aos candidatos em segundo turno que são removidos durante a etapa de
pré-processamento. Nesse sentido, a disparidade da variável é esperada, tendo em vista a própria
estrutura competitiva das eleições municipais, nas quais cada disputa comporta apenas um vencedor
por município. Esse desequilíbrio foi considerado nas etapas de treinamento e avaliação dos modelos,
conforme descrito nas subseções subsequentes.

No que tange à composição de gênero, observa-se que 1.883 candidatos são do sexo masculino (85%)
e 333 do sexo feminino (15%), distribuição que reflete a histórica sub-representação feminina nos
cargos executivos municipais brasileiros e que se mantém consistente com padrões identificados em
estudos anteriores sobre o perfil das candidaturas no país [12]. Entre as candidatas do sexo feminino,
a taxa de êxito eleitoral é proporcionalmente inferior à dos candidatos masculinos, o que pode levar
os modelos a carregar esse preconceito em suas predições [11].

\subsection{Pré-processamento}

\subsubsection{Limpeza e filtros}

Antes da modelagem, o conjunto de dados passou por uma etapa de limpeza orientada tanto pela
consistência lógica dos registros quanto pela prevenção de vazamento de informação. Registros
cuja situação de totalização constava como "Segundo Turno" foram removidos, pois candidatos nessa
condição aparecem duplicados na base, uma entrada referente ao primeiro turno e outra com o
resultado definitivo, o que introduziria redundância e potencial viés na estimação dos modelos.

Colunas de identificação direta, como o nome do candidato e a data de carga do registro, foram
descartadas por não carregarem valor preditivo e por representarem risco de memorização espúria.
O campo de município foi removido, mantendo apenas o seu respectivo código numérico, para melhor
compatibilidade com os modelos.

A variável de resposta foi recodificada a partir do campo de situação de totalização: candidatos
com o resultado "Eleito" receberam o valor 1, e os demais o valor 0, produzindo uma variável
binária diretamente interpretável pelos classificadores. Por fim, todas as variáveis categóricas
remanescentes — gênero, raça/cor, estado civil, faixa etária, grau de instrução, ocupação e partido
— foram convertidas para representações numéricas inteiras por meio de mapeamento ordinal, pelas
mesmas razões.

\subsubsection{Variável cor/raça}

*DEVE SER REVISADO*

O TSE não coletava autodeclaração racial em 2012. Os 649 candidatos daquele pleito (29,8% do
*dataset*) receberam o valor `#NE#` no campo `Cor/raça`, que foi mapeado para o grupo não-branco
na auditoria de *fairness*. Por isso, as métricas raciais comparam candidatos de raça branca
declarada com candidatos de raça não-branca ou com raça não registrada, não apenas com raça
não-branca.

Para estimar o efeito dessa ambiguidade, o experimento foi repetido excluindo 2012 (N = 1.529).
O DIR racial do XGBoost passou de 0,880 para 0,939 (ΔDIR = 0,059), acima do limiar de 0,05 usado
como critério de relevância substantiva. O TabPFN variou menos (DIR: 0,910 para 0,921; ΔDIR = 0,011).
A diferença indica que parte do desvio racial do XGBoost é associada à mistura temporal com os dados
de 2012, não unicamente ao atributo racial. O DIR < 1 do XGBoost deve ser lido como limite inferior
da disparidade real, e não como sua estimativa pontual.

\subsubsection{Codificação de variáveis}

*DEVE SER PREENCHIDO*

- Como cada variável categórica é representada nos modelos.  
- Diferença entre encoding para XGBoost e para TabPFN.

\subsubsection{Tratamento de desbalanceamento}

*DEVE SER REVISADO*

Dado o desbalanceamento entre classes descrito na Seção 3.1.4, foram adotadas estratégias distintas
para cada modelo, em função de suas características internas. No XGBoost, o parâmetro scale_pos_weight
foi ajustado pela razão entre o número de instâncias negativas e positivas, compensando diretamente
o desequilíbrio durante o treinamento. No TabPFN, o desbalanceamento é tratado implicitamente pelo
próprio mecanismo de inferência bayesiana do modelo, sem necessidade de ponderação explícita.

A calibração das probabilidades seguiu abordagens diferenciadas para cada classificador. Para o XGBoost,
a calibração foi realizada por regressão isotônica aplicada sobre previsões out-of-fold em validação
cruzada com cinco partições. Essa escolha se justifica pelo fato de que o XGBoost treinado com
scale_pos_weight tende a produzir escores polarizados, concentrados próximos aos extremos 0 e 1,
quando avaliado sobre o próprio conjunto de treino; calibração in-sample nesse cenário colapsa o
limiar ótimo τ* para a borda do grid de busca, comprometendo a utilidade das probabilidades estimadas.
Para o TabPFN, esse comportamento não se manifesta, de modo que a calibração foi conduzida in-sample
sem prejuízo à qualidade dos escores.

\subsection{Desenho experimental}

*TODA ESSA SUBSEÇÃO DEVE SER PREENCHIDA*

\subsubsection{Esquema de treino/teste}
- Split temporal entre eleições (treino em anos passados, teste no mais recente).  
- Detalhes: tamanhos de treino/teste e motivação temporal.

\subsubsection{Validação cruzada}  
- Se usar K-fold, como é aplicado (dentro do treino, estratificado etc.).  
- Número de folds, repetição ou não.

A proposta previa *Stratified 5-Fold Cross-Validation*. A implementação usou K = 10 (parâmetro
`NUMERO_DE_FOLDS` em `src/config.py`). Com K maior, cada fold de teste tem menos instâncias e a
variância da estimativa de MCC cai em torno de √2, o que favorece a comparação entre modelos.
O protocolo de 10 folds é padrão em benchmarks com conjuntos de dados na faixa de N ≈ 2.000 e
não invalida os testes estatísticos aplicados.

\subsubsection{Modelos avaliados}
- XGBoost baseline: hiperparâmetros, justificativa.  
- TabPFN: versão usada, forma de chamada, configuração.  
- XGBoost Blind: quais atributos sensíveis foram removidos (gênero e raça).

\subsubsection{Procedimento experimental}
- Passo a passo: para cada split, treinar três modelos, gerar predições, calcular métricas
de desempenho e fairness, salvar resultados.

\subsection{Métricas de avaliação}

*TODA ESSA SEÇÃO DEVE SER PREENCHIDA*

\subsubsection{Métricas de desempenho}  
- MCC (por que é adequado para desbalanceamento).  
- Outras (ex.: F1-macro, AUPRC se usar, accuracy para referência).

\subsubsection{Métricas de fairness}
- DIR (Disparate Impact Ratio): definição, valor ideal, threshold de 0,8.  
- EOD (Equal Opportunity Difference): definição, valor ideal.  
- Possível uso de FNR por grupo e predição positiva por grupo.

\subsubsection{Comparações estatísticas}  
- Teste de Wilcoxon sobre MCC entre modelos.  
- Tamanho de efeito que será reportado.

\section{Resultados}

\section{Conclusão}

\section{Referências Bibliográficas}

[1] Grari, V., Ruf, B., Lamprier, S., Detyniecki, M. Achieving Fairness with
Decision Trees: An Adversarial Approach. Data Sci. Eng. 5, 99–110, 2020.

[2] ROBERTSON, J. et al. FairPFN: Transformers Can do Counterfactual Fairness.
International Conference on Machine Learning (ICML), 2024.

[3] RAVICHANDRAN, S. et al. FairXGBoost: Fairness-aware Classification in
XGBoost. KDD Workshop on Machine Learning in Finance, 2020.

[4] GARG, A. et al. Real-TabPFN: Improving Tabular Foundation Models via
Continued Pre-training With Real-World Data. arXiv preprint, 2025.

[5] BASE DOS DADOS. Eleições Brasileiras. Dados do Tribunal Superior Eleitoral
(TSE) tratados. 2026.

[6] CHICCO, D.; JURMAN, G. The Matthews correlation coefficient (MCC) should
replace the ROC AUC as the standard metric for assessing binary classification.
BioData Mining, 2023.

[7] IMANI, M. et al. Why ROC-AUC Alone Is Insufficient for Highly Imbalanced
Data. Preprints.org, 2025.

[8] QU, J. et al. TabICL: A Tabular Foundation Model for In-Context Learning
on Large Data. arXiv preprint, 2025.

[9] WIKIPEDIA CONTRIBUTORS. TabPFN. Wikipedia, The Free Encyclopedia, 2025.
(Documentação sobre a versão v2.5 e funcionamento interno).

[10] referência data leakage

[11] MIGUEL, Luis Felipe e BIROLI, Flávia. Práticas de gênero e carreiras
políticas: vertentes explicativas. Rev. Estud. Fem. [online]. 2010, vol.18,
n.03, pp.653-679. ISSN 1806-9584.


\end{document}
